\section{Introduction}
\label{sec:intro}

Supercompilation~\cite{turchin-1986-supercompiler} fuses composed functions,
eliminates intermediate data structures, and discovers recursive structure
automatically.  Prior work established that standard supercompilation
corresponds to focused cyclic sequent proof search: driving is cut
elimination, generalisation is cut introduction, and the trace condition
is the global progress check of cyclic proof theory.\footnote{A companion
mechanisation (under review) formalises this correspondence in Rocq with
0 axioms; the present paper uses the resulting infrastructure as its
starting point.}  This paper shows that the correspondence is more than
a correctness proof --- it is a \emph{design instrument} that identified
two productive extensions to supercompilation.

\paragraph{Extension 1: LLM oracle for cut introduction.}
Syntactic anti-unification is a heuristic for cut introduction, but it is
purely structural and cannot discover semantic identities (map fusion,
accumulator introduction).  The sequent formulation suggests an
\emph{oracle architecture}: any function that proposes cut formulas is
valid if the kernel accepts the resulting proof graph.  We instantiate
this with an LLM (DeepSeek V4 Pro); the oracle is untrusted for soundness
--- the CIU theorem requires only the trace condition check.

\paragraph{Extension 2: LLM lemma proposal for the $\omega$-rule.}
Properties like $\mathsf{sorted}\ (\mathsf{sort}\ l) = \mathsf{true}$
require auxiliary lemmas.  The sequent formulation identifies this as
the $\omega$-rule: prove a lemma first, then use it as a rewrite rule.
We add a lemma environment to the SC: the LLM proposes auxiliary lemmas,
a sub-SC run proves them, and the main SC uses them as rewrite rules.
This enables compiler soundness and reverse-append distribution ---
properties unreachable by standard SC.

\paragraph{Extension 3: Dependent-type branch elimination.}
The SC's type annotations carry structural information.  We prune
impossible constructors from case-splits using type-index comparison
(e.g., \texttt{vnil} eliminated from $\mathsf{Vec}\ (\mathsf{succ}\ n)$).
This makes two programs with different dead branches produce identical
residuals, enabling more folding opportunities.

\paragraph{What is and is not claimed.}
The end-to-end pipeline is demonstrated: the LLM proposes a lemma,
the sub-SC validates it, and the SC uses it as a rewrite rule.  The scrambling experiment
(Section~\ref{sec:eval}) uses 6 programs and 24 tests, finding that
the LLM succeeds with high confidence on all 6 even with all function
names scrambled.  The type-index trace condition
extension (\S\ref{sec:design}) is specified and implemented.

\paragraph{Contributions.}
\begin{enumerate}[leftmargin=*]
\item \textbf{LLM oracle architecture} (Section~\ref{sec:llm}).
  A 2-layer generalisation strategy: anti-unification $\to$ LLM oracle.
  All LLM proposals are validated by the kernel.  The scrambling experiment
  shows the LLM succeeds on 6 programs even with all names scrambled,
  and with definitions provided, all 6 pass with high confidence.

\item \textbf{Mechanised $\omega$-rule} (Section~\ref{sec:limits}).
  A lemma environment: lemmas are proved by sub-SC and used as rewrite
  rules in the main SC.  This enables compiler correctness and
  reverse-append distribution.  The LLM proposal step succeeds on
  all tested examples (Section~\ref{sec:eval}).

\item \textbf{Dependent-type branch elimination} (Section~\ref{sec:design}).
  Type-index comparison prunes impossible constructors during case-splitting.
  The pruning result shows two programs with different dead branches produce
  identical residuals --- a strictly smaller output than unpruned SC.
\end{enumerate}

\paragraph{Empirical results.}
65 theorems proved by \vmcompute{}.  The breakdown:\footnote{Excludes
infrastructure lemmas (CanonRoundtrip.v, 16 theorems) from the primary
count.  Including them: 81 theorems.  All 0 Admitted.}
\begin{itemize}[leftmargin=*]
\item \textbf{Standard SC} (AU only): 33 theorems --- monad laws, map fusion,
  sort-length, sort-sum, insert-length, take-drop.
\item \textbf{$\omega$-rule} (lemma-driven): 16 theorems --- reverse-append
  distribution, sort correctness, compiler soundness, gradual cast safety.
\item \textbf{Dependent-type pruning}: 4 theorems --- Vec index
  normalisation, dead-branch elimination.
\item \textbf{LLM evaluation}: 12 theorems --- scrambling experiment.
\end{itemize}

\paragraph{Outline.}
Section~\ref{sec:background} reviews the SC-as-proof-search dictionary.
Section~\ref{sec:design} explains how the sequent perspective identified
the gaps.  Sections~\ref{sec:llm} and \ref{sec:limits} present the
LLM oracle and lemma proposal.  Section~\ref{sec:examples} presents
the full example suite including the dependent-type results.
Section~\ref{sec:eval} reports the scrambling experiment.
Section~\ref{sec:related} positions the work.
